%-*-latex-*-
%-*-latex-*-
\newcommand\COURSE{ciss240}
\newcommand\ASSESSMENT{q1007}
\newcommand\ASSESSMENTTYPE{Quiz}
\newcommand\POINTS{\textwhite{xxx/xxx}}

\input{myquizpreamble}
\input{yliow}
\input{\COURSE}
%-*-latex-*-

%-*-latex-*-

%-*-latex-*-
\renewcommand\TITLE{\ASSESSMENTTYPE \ \ASSESSMENT}

\renewcommand\EMAIL{}
\renewcommand\AUTHOR{eagraham1@cougars.ccis.edu}

\textwidth=6in
\begin{document}
\topmatterthree

This is a closed-book, no compiler, 5 minute quiz.

%------------------------------------------------------------------------------
\nextq
The following is a plot of
\[
f(x) = 2 x^{2.3} \cos{(x^{4.5} + 1)} + 20
\]
\begin{python}
from latextool_basic import *
p = FunctionPlot(domain='0:3.5', height='1.5in')
p.add('2 * x**2.3 * cos(x**4.5 + 1) + 20', color='red', python=1, pin='below', num_points=2000,
pin_message='')
print(p)
\end{python}
Write a program that prints the value of $x$ from $a=1$ to $b=1.5$ where
$f(x)$ reaches a maximum value using a  (brute force) search with
1000 equally spaced points on the interval $[1, 1.5]$.
The printout must be in fixed point format with 5 decimal places.
Skeleton code is provided.
\\
\ANSWER
\begin{answercode}
#include <iostream>
#include <iomanip>
#include <cmath>

int main()
{
    double max_x = 0.0;               // x where the max of the function occurs
    double max_y = -1000000000.0;     // max value of functions
    double a = 1.0, b = 1.5;          // endpoints
    int n = 1000;                     // number of points
    double dx = (b - a) / n;          // distance between consecutive points

    for (int i = 0; i <= n; ++i)      
    {
        double x = a + i * dx;        
        double y = 2 * pow(x, 2.3) * cos(pow(x, 4.5) + 1) + 20;)

        if (y > max_y)
        {
            max_y = y;                // Update max_y
            max_x = x;                // Update max_x
        }
    }
    std::cout << std::fixed << std::setprecision(5)
              << "Max x: " << max_x << std::endl
              << "Max y: " << max_y << std::endl;

    return 0;
}
\end{answercode}

%------------------------------------------------------------------------------
\nextq
Write a program that gets a positive integer $n$ from the user
and prints the sum of the digits of $n$.
For instance if the user enters \verb!2407!, the program prints
\verb!13!:
\begin{Verbatim}[frame=single,fontsize=\footnotesize,commandchars=\\\{\}]
\userinput{2407}
13
\end{Verbatim}
Skeleton code is provided.
\\
\ANSWER
\begin{answercode}
#include <iostream>
#include <cmath>

int main()
{
    int n;
    std::cin >> n;

    // compute the number of digits in n and store in num_digits
    int num_digits;

    // compute the sum of digits in n and store in sum_digits
    int sum_digits = 0;

    while(n > 0)
    {
         sum_digits += n % 10;
         n /= 10;
         num_digits++;
    }
    
    std::cout << sum_digits << '\n';
    return 0;
}
\end{answercode}

%------------------------------------------------------------------------------
%-*-latex-*-
\newpage
\input{instructions}
\end{document}

